\section{总结与收获}

本课程设计围绕云计算平台搭建与 Spark 分布式数据分析两个方向展开，涵盖了 Docker 容器化、Kubernetes 集群部署、配置管理、持久化存储、弹性伸缩以及 Spark 大数据处理等多个核心内容。通过在华为云真实环境中的实践操作，不仅完成了任务书要求的各项实验内容，也进一步加深了对云原生架构和分布式计算体系的理解。本章结合项目实施过程，对课程设计中的主要收获、技术体会以及后续改进方向进行总结。

\subsection{云计算平台搭建总结}

\subsubsection{容器化与Kubernetes部署理解}

在传统软件部署模式下，应用通常依赖固定服务器环境，不同系统之间容易出现依赖版本不一致、运行环境差异等问题。而在本课程设计中，通过 Docker 技术将应用及其运行环境封装为标准镜像，实现了“一次构建，到处运行”的部署模式。

在容器化实践过程中，进一步理解了镜像、容器、镜像仓库以及容器编排之间的关系。Dockerfile 负责定义应用运行环境，Docker Compose 用于本地多容器协同验证，而 SWR 则承担镜像集中存储与分发功能。通过这一过程，认识到容器技术不仅能够简化部署流程，还能够显著提高应用迁移与扩展能力。

在 Kubernetes 部署阶段，重点学习了 Deployment、Pod 和 Service 等核心资源对象的作用。其中 Deployment 负责应用副本管理与滚动更新，Pod 是 Kubernetes 调度的最小运行单元，而 Service 则为 Pod 提供统一访问入口和服务发现能力。通过实际部署后端、前端和 Redis 服务，进一步理解了 Kubernetes 如何实现应用生命周期管理以及故障自动恢复机制。

此外，通过对 ELB、LoadBalancer Service 和公网访问流程的实践，也更加清晰地认识到云平台网络架构中负载均衡、服务暴露以及流量转发的实现方式。相比本地实验环境，真实云平台更能够体现现代云原生应用的部署模式和运行机制。

\subsubsection{配置分离、持久化与弹性伸缩理解}

在项目实施过程中，ConfigMap、Secret、PVC 和 HPA 等 Kubernetes 高级资源的使用，使配置管理和运维能力得到了进一步提升。

首先，ConfigMap 与 Secret 的应用使系统配置与应用代码实现了解耦。Redis 地址、端口等非敏感配置通过 ConfigMap 管理，而密码等敏感信息则通过 Secret 存储。这种方式避免了将环境配置硬编码到程序中，提高了系统的可维护性和安全性。

其次，通过 Redis 持久化实验，对 PVC 和持久化存储机制有了更加深入的理解。在实验中，即使删除 Redis Pod，数据仍能够在新创建的 Pod 中继续访问，说明 PVC 已经成功将数据与容器生命周期分离。该过程体现了云原生架构中“计算资源可销毁、数据资源可保留”的设计思想。

最后，通过 HPA 自动伸缩实验，进一步认识到云平台资源动态调度的重要性。当系统负载升高时，HPA 能够自动增加应用副本，提高服务处理能力；当负载降低时，又能够自动减少副本数量，从而节约资源成本。这种弹性伸缩能力是云计算区别于传统服务器部署的重要优势之一，也是实现资源高效利用的重要手段。

\subsection{Spark大数据分析总结}

\subsubsection{分布式数据处理流程理解}

在 Spark 大数据分析部分，项目以豆瓣电影数据集为对象，完成了数据清洗、统计分析以及性能测试等任务。通过实际编写 PySpark 程序，对 Spark 的分布式执行机制有了更加直观的认识。

在 Spark 架构中，Driver 负责作业解析、任务调度和执行控制，而 Executor 负责实际的数据计算。用户提交的程序首先由 Driver 转换为执行计划，再将任务分发给多个 Executor 并行执行，最终汇总计算结果。与传统单机程序相比，Spark 能够充分利用集群资源处理更大规模的数据集。

在数据清洗过程中，学习了 Spark DataFrame 的基本使用方法，包括数据加载、Schema 推断、缺失值处理以及统计分析等操作。通过 DataFrame API，可以方便地完成数据转换和预处理任务，并且能够自动优化执行计划，提高运行效率。

在 Spark SQL 部分，进一步掌握了 GROUP BY 聚合分析、Top-N 排序分析、时间趋势分析以及窗口函数分析等常见数据分析方法。与传统数据库 SQL 不同，Spark SQL 可以直接作用于分布式数据集，从而兼顾 SQL 语言的易用性和分布式计算的高性能特点。

通过 Spark Operator 在 Kubernetes 环境中运行 SparkApplication，也进一步理解了 Spark 与云原生平台结合的实现方式。相比传统独立集群部署模式，Spark on Kubernetes 能够利用 Kubernetes 提供的资源调度和弹性管理能力，实现更加灵活的分布式计算环境。

\subsubsection{性能对比与Amdahl定律理解}

在性能分析实验中，通过比较 Pandas、单 Executor PySpark 和双 Executor PySpark 的执行时间，进一步理解了并行计算的性能特征。

实验结果表明，增加 Executor 数量并不一定能够获得完全线性的性能提升。这一现象与 Amdahl 定律的结论一致。根据 Amdahl 定律，程序整体加速比不仅取决于并行部分的执行效率，还受到串行部分比例的限制。当程序中仍存在大量无法并行化的工作时，即使增加计算资源，整体性能提升也会逐渐趋于饱和。

在 Spark 实际运行过程中，任务启动、资源调度、数据序列化、Shuffle 通信以及结果汇总等操作均属于额外开销。当数据规模较小时，这些开销甚至可能超过计算本身的时间消耗。因此，在小规模数据集上，PySpark 的优势并不会十分明显；而在大规模数据场景下，分布式计算的性能优势才会逐渐体现出来。

通过将理论模型与实验结果进行对比分析，不仅验证了 Amdahl 定律的实际意义，也进一步认识到分布式系统设计需要综合考虑计算规模、通信开销和资源利用率等因素，而不能简单依赖增加节点数量提升性能。

\subsection{课程设计反思}

\subsubsection{主要挑战}

本课程设计最大的挑战来自于真实云平台环境的复杂性。与课堂实验相比，云环境涉及网络配置、镜像仓库、负载均衡、存储资源以及 Kubernetes 多种资源对象之间的协同工作，任何一个环节出现配置错误都可能导致整体部署失败。

此外，Spark 分布式分析部分涉及镜像构建、依赖安装、数据集管理以及 Kubernetes 作业调度等多个方面，问题定位过程较为复杂。例如分析脚本缺失、数据文件缺失、镜像缓存未更新以及 Python 依赖库缺失等问题，都需要结合日志信息逐步排查才能发现根本原因。

通过解决这些问题，不仅提高了实际工程排错能力，也更加深刻地理解了云计算平台中各组件之间的协作关系。

\subsubsection{改进方向}

虽然本项目已经完成了课程设计要求的主要功能，但仍存在进一步完善的空间。

首先，可以引入更加规范的 CI/CD 流水线，实现镜像自动构建、自动测试和自动部署，进一步提升项目交付效率。其次，可以将 Spark 数据源扩展到 OBS 对象存储，从而更接近真实大数据处理场景。再次，可以增加监控与日志系统，例如 Prometheus、Grafana 和 ELK，实现应用运行状态和资源使用情况的实时可视化。

在性能分析方面，也可以采用更大规模的数据集和更多 Executor 配置进行测试，以获得更加具有代表性的性能对比结果，并进一步验证分布式计算模型在不同规模场景下的表现差异。

总体而言，本课程设计不仅完成了课程要求的实验任务，更重要的是通过真实云平台实践，将课堂中学习的云计算、容器技术、Kubernetes 和分布式计算知识进行了系统整合与应用，为后续学习和工程实践奠定了良好基础。
